% Owner: Duong. Claims are bounded to what the cited studies report; see the
% source-to-claim notes in the project research artifacts.
\section{Related Work}\label{sec:related}

\subsection{Traditional Rental Management Systems}
A rental platform could be built as one application backed by one database for listings, requests, contracts, payments, and maintenance records. Such a design is a candidate comparison, not a measured baseline in this paper. It could place changes that must commit together in one local transaction and use database joins for cross-domain reports. Whether maintenance and billing require different release cycles or scaling profiles is a workload question for the rental system, not a consequence established by this paper.

Table~\ref{tab:related} compares the specific studies cited in this paper. Those studies address microservice boundaries, data management, event coordination, and performance in other contexts; they do not evaluate the proposed rental workflow. This selective comparison is not a systematic review of rental-management systems. A centralized rental implementation would need to be built and tested under matched conditions before serving as an empirical baseline.

\subsection{Microservices Architecture}
Microservice decomposition assigns a business capability, and its data, to a service boundary. The research base is broad: a systematic mapping classified 103 primary studies and assessed their industrial-adoption potential \cite{dipfrancesco2019architecting}. That map shows a mature body of design knowledge rather than a settled answer for any domain, and it derives no boundaries for rental management.

Boundary identification is where empirical work is most instructive. A review of 34 studies plus a survey of 63 practitioners found wide variation in understanding of domain-modeling artifacts and continuing need for methodological support \cite{zhong2024ddd}. Domain-driven design is thus a reasoned lens for decomposing a rental workflow, not a procedure yielding an objectively optimal partition. A mixed-method study of microservice data management reaches a compatible conclusion, identifying ownership, fault and performance isolation, and independent schema evolution as recurring practitioner concerns \cite{laigner2021datamanagement}. The database-per-service rule in Section~\ref{sec:architecture} follows from that ownership concern, not from evidence that a seven-service partition is optimal.

\subsection{Event-Driven Architecture}
Events record a completed local fact and let other services react without a synchronous command chain. The coordination semantics predate microservices: a Saga decomposes a long-lived transaction into interleavable steps completed either by all steps succeeding or by compensating the completed ones \cite{garciamolina1987sagas}. Comparing four Java frameworks on an order workflow, \v{S}tefanko et al.\ show that committing each step separately makes partial results visible, relaxing isolation and moving the system to eventual consistency, where compensation semantically undoes work rather than restoring the prior state \cite{stefanko2019saga}. They also observed duplicate orders under rapid requests in one implementation---a test observation, not a property of the pattern, but illustrative of why consumer-side deduplication cannot be assumed. A log-based broker retains the record durably, yet does not by itself establish that each consumer applied it \cite{kleppmann2017designing}.

Operational cost is documented independently of performance claims. Mining more than 8,000 developer questions, Laigner et al.\ identified large payloads, schema modeling, flow debugging, and dependency and ordering constraints as recurring difficulties \cite{laigner2026eventmanagement}. In a comparison of one application implemented as a monolith and as an event-driven system, the monolith consumed fewer resources and had better response times \cite{cabane2024edaperformance}. This single-application result does not establish a general ranking; it motivates matched workloads, a stated environment, and measured outcomes here.

Two proposed mechanisms differ in evidential status. The transactional outbox appears in an applied redesign moving gateway calls to Kafka events, and as a Spring methodology for publishing within a database transaction \cite{peliza2025outbox,glybovets2024transactions}; neither reports a controlled reliability comparison, and the first calls its validation partial. Idempotence is well studied at service level as a route to fault tolerance \cite{ramalingam2013idempotence}, but does not address broker-side consumer deduplication, which is what a consumer relies on.

\subsection{Existing Approaches and Research Gaps}
The cited studies do not supply an evaluated service partition or fault results for this particular rental workflow. Section~\ref{sec:architecture} therefore argues for its boundaries from the versioned requirements, while Section~\ref{sec:evaluation} defines tests for acceptance-to-draft delivery, duplicate effects, and recovery. The protocol is intended to characterize this implementation; it cannot establish superiority over a centralized design without a matched comparison.

% Owner: Duong. This is a comparison of source scope, not a claim of measured superiority.
\begin{table*}[t]
\caption{Related Work Comparison and Evidence Needed for the Rental Case}
\label{tab:related}
\centering\footnotesize
\begin{tabularx}{\textwidth}{@{}p{0.19\textwidth}X X X@{}}
\toprule
Source or approach & Scope established by source & Relevance to this design & Evidence still needed \\
\midrule
Richardson \cite{richardson2018microservices} & Microservice decomposition and integration patterns & Frames service ownership and outbox-style coordination & Rental-specific implementation and comparative outcomes \\
Kleppmann \cite{kleppmann2017designing} & Distributed data, logs, and failure semantics & Frames delivery, replay, and consistency questions & Workload-specific measurements and failure traces \\
Traditional centralized design & Candidate single-application baseline; source pending & Can provide local transaction and deployment comparison & Verified rental-system study and reproducible baseline \\
This paper's proposed design & Seven named rental services and a planned evaluation & Makes domain boundaries and failure experiments explicit & Specification audit, implementation evidence, and results \\
\bottomrule
\end{tabularx}
\end{table*}

